% 家庭.tex
% Auto-generated from 感悟论语.odt by tools/odt_to_tex.py

\setpartplant{artwork/xingcai.png}
  {荇菜 · 《诗·周南·关雎》}
  {父母在，不远游，游必有方。}
  {《论语·里仁》}
  {家庭}
\part{家庭}

\chapter{色难}

“色难”要求我们在从内心和行动上对于父母实践孝的内容外，还应始终保持足够的耐心与和颜悦色的态度。它看似极其简单，实则很难做到。

通常，我们的逐渐成熟与父母的逐渐衰老是一刻不停地持续发生的并行过程——也正因为如此，才会有“子曰：‘父母之年，不可不知也，一则以喜，一则以惧。’（《论语 里仁篇第四 4.21》）”。所以，通常当我们自己与时俱进的时候，父母则不可避免地有所落伍于时代。与此同时，父母长期以来积累形成的习惯和观点也都在一天天地不断固化——事实上，我们中的大多数也经历着同样的过程，只不过相对于父母有所延迟罢了。这样一来，我们在观点看法、做事方式、生活习惯等诸多方面就难免会与父母之间有不一致的地方，甚至还有争执，尽管大都也只是就事论事的。

再者，即便是我们基于科学理性的思考和调查而给父母提出的良好、合理化建议，在父母那里能够得到认可、接受或实践却变得很困难。比如，良好的饮食、起居和锻炼习惯，健康豁达的心态等。哪怕我们能够在一开始就耐下心来给父母苦口婆心地讲道理或做科普，也很有可能因为父母的接受能力和思想与心态的开放程度而不见成效。

这个时候为了做到不给父母脸色看，就需要想方设法用父母能够理解和接受的方式去讲解、交流和沟通。这包括在技术方面可以使用先进的可视化手段，在语言上则不妨试试具体和形象的比喻或类比。可以想到，循循善诱、传道解惑的老师所应具备的能力和修养，正是我们需要借鉴和采纳的。

\begin{classic}
《论语 为政篇第二 2.8》子夏问孝。子曰：“色难。有事，弟子服其劳，有酒食，先生馔，曾是以为孝乎？”
\end{classic}

\entrydate{2018.02.22}

\chapter{养志胜于养口体}

昨日收到母亲买的土法工艺制成的豆瓣酱一罐。罐体制作精良，打开后里面的酱也非常香。以前我自己买过李锦记等牌子，但都没有这次的好。

由此想到，父母虽已退休多年，体力与精力也比不上年轻的时候。但他们想要为儿女分担家庭生活中的琐事、在力所能及的事情上想要出力搭把手的心思始终未变。更何况，我目前居于外地，而他们仍在老家，相隔千里，他们能够时时刻刻为我做事情的机会就更少了。因此，要发现、抓住、创造机会，让他们能够为我平日的学习、工作、生活出一分力。这样一来，父母便会觉得他们自己还能为子女做很多事、帮很多忙、解决很多问题，自然地在心态上就会始终保持年轻。

同理，每次我假期回家后返程时，父母总会让我带上很多家里自制的食品、家乡的特产。相较于年少时的不更事，现在的我总是欣然接受，不会嫌行李太重而拒绝。这其中的思想源于《孟子·离娄上》：

\begin{classic}
曾子养曾晳，必有酒肉。将撤，必请所与。问有余，必曰“有”。曾晳死，曾元养曾子，必有酒肉。将撤，不请所与。问有余，曰“亡矣”。将以复进也。此所谓养口体者也。若曾子，则可谓养志也。
\end{classic}

所以，对于父母应养其志，而不只是口体。

相反地，我们不能认为上述做法会给父母添麻烦。实际上，人与人之间缘分和感情的维系，就是靠着一系列“麻烦”的互动而得以巩固、坚实，以至弥足珍贵。

\entrydate{2023.07.31}

\chapter{为什么要与父母讨论自己的科研工作}

为什么要与父母讨论自己的科研工作，尽管他们其实并不懂我的专业？这其实是儒家思想所要求的：不能只养父母的口体，还要养心志。

\begin{classic}
《孟子·离娄上》：曾子养曾皙，必有酒肉。将彻，必请所与。问有余，必曰：‘有。’曾皙死，曾元养曾子，必有酒肉。将彻，不请所与。问有余，曰：‘亡矣。’——将以复进也。此所谓养口体者也。若曾子，则可谓养志也。事亲若曾子者，可也。”
\end{classic}

父母生于上世纪50年代，受文革影响并未念太多书并接受完整、良好的中小学教育。以前我上小学的时候，他们指导我做功课的情形依旧历历在目。如今我通过向他们讲述我的科研工作与算法开发项目的进展，听取他们的宏观性与原则性的建议，就能让他们重温指导我学习的记忆、拥有参与到我自己事业中的感觉，并也能分享我取得每一点点进步时的喜悦与成就感。这样一来，他们便会觉得依然能够发挥作用、产生价值，从而保持年轻的心态。

\entrydate{2025.05.08}

\chapter{不知老之将至云尔}

去年五一回家期间给父母用上了树莓派，并将之前已淘汰的台式机中的 500 GB 硬盘拆了下来接在上面，作为数据仓库。虽然该硬盘只有 5000 rpm 的转速，且缓存仅有 16 MB ，所以数据传输速率欠佳，但这对于父母日常浏览照片、查看歌谱倒也基本够用了。然后，通过树莓派上的 Samba 文件服务器，让他们可以从笔记本电脑和 iPad 远程访问这些数据。

今年春节回家，我给父母配置好并教会了 Copilot 的使用。争取每年都能介绍和引入一门新技术，与时俱进。正所谓“发愤忘食，乐以忘忧，不知老之将至云尔”。

\begin{classic}
《论语·述而篇第七 7.19》叶公问孔子于子路，子路不对。子曰：「女奚不曰，其为人也，发愤忘食，乐以忘忧，不知老之将至云尔。」
\end{classic}

同时，论语中还有

\begin{classic}
《论语·里仁篇第四 4.21》子曰：「父母之年，不可不知也：一则以喜，一则以惧。」
\end{classic}

以及

\begin{classic}
《论语·颜渊篇第十二 12.4》司马牛问君子。子曰：「君子不忧不惧。」曰：「不忧不惧，斯谓之君子矣乎？」子曰：「内省不疚，夫何忧何惧？」
\end{classic}

所以，正是通过帮助父母与时俱进，让他们“忘老”，永葆年青，才能做到君子不惧。

\entrydate{2024.02.18}
